\section{Discussion and Limitations}

Projection--Carry Geometry is not a new multiplication algorithm. It is a
structural and geometric framework for the ordinary base-\(B\) multiplication
process. The goal is to separate bilinear projection from nonlinear carry flow
and to make the geometry of carry propagation measurable.

Support geometry controls potential mass concentration, not exact carry. The
sumset \(S_X+S_Y\), representation counts \(\rho_{X,Y}(k)\), and additive
energy describe where digit interactions can collide. They do not determine
the actual raw coefficient values. Digit amplitudes determine \(c_k\), and
incoming carry determines the realized activation condition \(c_k+r_k\ge B\).
Thus carry is locally activated but globally propagated through the recurrence.

The experiments are synthetic random digit experiments. They are useful for
testing structural hypotheses and producing reproducible scaling tables, but
they do not characterize every input distribution that appears in arithmetic
software or hardware. Deterministic support families may display different
behavior.

Normalized metrics must also be interpreted carefully. Normalization by digit
length, support size, or total interaction count can remove the scale that
drives absolute carry mass. For this reason, normalized figures should be read
alongside raw carry and raw support-concentration plots.

Future work includes binary multiplication and bit-level circuits,
multi-factor tensor projection, signed digits, negative coefficients, carry
minimization, and more direct connections to additive combinatorics.
